\documentclass[11pt]{article}

% ============================================================================
% periods_quasiperiods_determinant.tex
%
% Standalone explanatory note (NOT manuscript text): where the
% period/quasi-period determinant relation comes from, and why "period"
% / "quasi-period" / "first kind" / "second kind" are the right words.
% Side-by-side: elliptic curve (dx/y, x dx/y) vs. modular (Delta, Dhat).
%
% Sourced from the project knowledge:
%  - Brown, "Quasi-periods of some cusp forms..." (QP.pdf): de Rham
%    construction, period matrix (eta^+ omega^+ ; i eta^- i omega^-),
%    det proportional to a power of 2 pi i.
%  - Brown 1710.07912v2 §1.3, §2: M^!_{n+2}/D^{n+1}M^!_{-n} = H^1_dR,
%    "modular forms of the second kind", Bol operator D = q d/dq,
%    period matrix P_f, det(P_f) proportional to power of 2 pi i.
%  - Fonseca, "On coefficients of Poincare series..." : explicit EC
%    example (omega,eta)=(dx/(2y+1), x dx/(2y+1)), det P = 2 pi i
%    (Legendre), and the modular dictionary X = sigma_1^v, Y = -sigma_tau^v,
%    comp(omega) = 2 pi i (X - tau Y).
%  - Zagier 1991 (MF_Periods_and_Jacobi_theta): k=12 values, the rank-1
%    Lemma forcing omega^pm / eigenform alignment, and
%    omega^+ omega^- / i(Delta,Delta) = 2^10.
%
% Standalone; compiles alone; no \ref/\cite into the main file.
% ============================================================================

\usepackage[margin=1.05in]{geometry}
\usepackage{amsmath, amssymb, amsthm}
\usepackage{mathtools}
\usepackage{array}
\usepackage{xcolor}
\usepackage{hyperref}
\hypersetup{colorlinks=true, linkcolor=blue!55!black, urlcolor=blue!55!black}

\newcommand{\HH}{\mathbb{H}}
\newcommand{\ZZ}{\mathbb{Z}}
\newcommand{\QQ}{\mathbb{Q}}
\newcommand{\CC}{\mathbb{C}}
\newcommand{\RR}{\mathbb{R}}
\newcommand{\Dhat}{\widehat{\Delta}}
\newcommand{\HdR}{H^1_{\mathrm{dR}}}
\newcommand{\HB}{H^1_{\mathrm{B}}}
\DeclareMathOperator{\Sym}{Sym}

\newtheorem*{slogan}{In one line}

\title{\bf Periods, quasi-periods, and where the determinant comes from\\[2pt]
\large An elliptic-curve / modular-form Rosetta stone}
\author{Working note --- conceptual, not manuscript text}
\date{\today}

\begin{document}
\maketitle

\noindent\emph{Goal: explain why $\omega^+\eta^- - \omega^-\eta^+$ is a
power of $2\pi i$ (up to rationals), and why Brown's words --- ``period,''
``quasi-period,'' ``first kind,'' ``second kind'' --- are the sensible
ones. The whole story is one structural fact (a $2\times 2$ period matrix
between a de Rham and a Betti basis, with a symplectic pairing) seen
twice: once for an elliptic curve, where it is the classical Legendre
relation, and once for $\Delta$, where it is Brown's determinant. The
determinant is an intersection number, and that is why your construction
should be able to \emph{derive} it rather than merely match it
numerically.}

\bigskip\hrule\bigskip

\section{The elliptic curve, concretely (the case to anchor on)}

Let $E$ be an elliptic curve over $\QQ$. Its first algebraic de Rham
cohomology $\HdR(E)$ is $2$-dimensional, and it has a distinguished basis
of \emph{differentials}:
\[
  \omega = \frac{dx}{2y+1} \quad(\text{first kind}),
  \qquad
  \eta = \frac{x\,dx}{2y+1} \quad(\text{second kind}).
\]
``First kind'' means \emph{holomorphic} (no poles anywhere, including at
infinity); it is the everywhere-regular differential. ``Second kind''
means it has poles but \emph{no residues} --- so it still defines a
cohomology class (a residue would make the period path-dependent /
ill-defined). $\eta$ has a double pole at the origin of $E$ and no
residue. The pair $(\omega, \eta)$ is a basis of $\HdR$, and the second
basis vector is genuinely \emph{not} holomorphic --- that is the entire
point of needing two of them.

On the Betti side, $H_1(E(\CC); \ZZ) \cong \ZZ^2$ has a basis of
\emph{cycles} $(\gamma_1, \gamma_2)$ (the two loops of the torus). Pairing
differentials against cycles --- integrate a form over a loop --- gives
the \textbf{period matrix}
\[
  P = \begin{pmatrix}
    \int_{\gamma_1}\omega & \int_{\gamma_1}\eta \\[3pt]
    \int_{\gamma_2}\omega & \int_{\gamma_2}\eta
  \end{pmatrix}
  = \begin{pmatrix} \omega_1 & \eta_1 \\ \omega_2 & \eta_2 \end{pmatrix}.
\]
The left column ($\omega_1, \omega_2$) are the \textbf{periods} --- the
integrals of the holomorphic differential, the lattice generators. The
right column ($\eta_1, \eta_2$) are the \textbf{quasi-periods} --- the
integrals of the second-kind differential. (Classically, with
$\omega = dz$ and $\eta = \wp(z)\,dz$ in Weierstrass coordinates, the
$\eta_i$ are the values of the Weierstrass $\zeta$-function increments;
``quasi-period'' is the period of the second-kind partner of the period.)

\begin{slogan}
Periods integrate the first-kind (holomorphic) differential; quasi-periods
integrate its second-kind (meromorphic, residue-free) partner. Both are
columns of one $2\times 2$ period matrix.
\end{slogan}

\section{Legendre's relation \emph{is} the determinant}

The determinant of $P$ is not a generic transcendental --- it is forced:
\[
  \boxed{\ \det P
  = \omega_1\eta_2 - \omega_2\eta_1
  = 2\pi i\ }\qquad(\text{Legendre's period relation}).
\]
(With the normalization $(\omega,\eta) = (dx/(2y+1),\, x\,dx/(2y+1))$ this
is exactly $2\pi i$; in Fonseca's worked $\QQ$-curve it is verified
numerically to $\det P = 2\pi i$ on the nose.)

Why is the determinant special, when the individual entries are
honest transcendentals (the $\omega_i$ alone are not algebraic)? Because
$\HdR(E)$ carries a \textbf{symplectic pairing} --- the cup product / the
algebraic intersection pairing on $H^1$ --- and $\langle \omega, \eta
\rangle$ is a \emph{rational} (here, essentially integer) number, since it
is computed algebraically from the differentials. The Betti side carries
the dual pairing, the topological intersection number $\gamma_1 \cdot
\gamma_2 = 1$ of the two torus loops. The comparison isomorphism
$\HdR \otimes \CC \cong \HB \otimes \CC$ --- ``integrate the form over the
cycle,'' i.e.\ the matrix $P$ --- must \emph{intertwine} the two pairings,
and the factor by which a comparison isomorphism scales a symplectic form
is its determinant. Concretely:
\[
  \underbrace{\langle \omega, \eta\rangle_{\mathrm{dR}}}_{\in\,\QQ}
  \;=\; \frac{1}{2\pi i}\,
  \det P \cdot \underbrace{(\gamma_1\cdot\gamma_2)}_{=\,1}.
\]
So $\det P$ is pinned to $2\pi i$ times a rational \emph{because} it is the
ratio of two rational symplectic volumes (de Rham and Betti), with the
$2\pi i$ the normalization of the comparison on $H^1$ of a curve. The
determinant is an \textbf{intersection number in disguise}; that is the
whole content of Legendre.

\begin{slogan}
$\det P = 2\pi i$ holds because $P$ matches a rational de Rham symplectic
pairing to a rational Betti one; the determinant is the scaling factor,
and $2\pi i$ is the normalization of the comparison on a curve. Legendre =
``the periods respect the intersection form.''
\end{slogan}

\section{The modular form, same skeleton}

Now replace the curve by the \emph{family} of all elliptic curves --- the
modular curve / moduli stack $\mathcal{M}_{1,1}$ --- and take symmetric
powers of the rank-$2$ bundle. Brown's Theorem (1710.07912 §2; QP §1) is
the family version of ``$\HdR(E)$ is $2$-dimensional with a first- and
second-kind basis'':
\[
  M^!_{n+2}(\QQ)\big/ D^{\,n+1} M^!_{-n}(\QQ)
  \;\xrightarrow{\ \sim\ }\;
  \HdR\!\big(\mathcal{M}_{1,1};\, \Sym^n\big),
  \qquad D = q\frac{d}{dq}.
\]
Read this slowly, because it is the crux. The left side is weakly
holomorphic modular forms of weight $n+2$, \emph{modulo} $(n{+}1)$-fold
$q$-derivatives of weight-$(-n)$ forms (Bol's operator $D^{n+1}$ raises
weight $-n$ to $n+2$). The quotient is what Brown calls \textbf{modular
forms of the second kind}, and the theorem says it \emph{is} the de Rham
cohomology of the local system. So:
\begin{center}
\begin{tabular}{>{\raggedright}p{0.46\textwidth} >{\raggedright\arraybackslash}p{0.46\textwidth}}
\textbf{Elliptic curve} & \textbf{Modular ($k=12$, $n=10$)}\\
\hline
$\HdR(E)$, rank $2$ &
$\HdR(\mathcal{M}_{1,1}; \Sym^{10})$, rank $2$ per Hecke eigenspace\\[4pt]
$\omega = dx/(2y{+}1)$, first kind (holomorphic) &
$\Delta \in S_{12}$, the holomorphic cusp form\\[4pt]
$\eta = x\,dx/(2y{+}1)$, second kind (pole, no residue) &
$\Dhat \in S^!_{12}$, weakly holomorphic --- the second-kind partner\\[4pt]
de Rham symplectic pairing $\langle\omega,\eta\rangle \in \QQ$ &
the Poincar\'e / cup pairing on $\HdR(\Sym^{10})$, $\QQ$-valued\\[4pt]
Betti cycles $\gamma_1, \gamma_2$, $\gamma_1\cdot\gamma_2 = 1$ &
the cuspidal cohomology $H^1_{\mathrm{cusp}}(\Gamma; V_{10})$, rank $2$\\[4pt]
period matrix $P = \big(\int_{\gamma_i}\omega \ \big|\ \int_{\gamma_i}\eta\big)$ &
$P_f = \begin{psmallmatrix} \eta^+ & \omega^+ \\ i\eta^- & i\omega^- \end{psmallmatrix}$ (Brown)\\[8pt]
$\det P = 2\pi i$ (Legendre) &
$\det P_f \doteq (2\pi i)^{k-1}$ (up to $\QQ^\times$)\\
\end{tabular}
\end{center}

The dictionary is not analogy-by-vibes; it is literally the same theorem
applied to $\Sym^n$ of the curve's $H^1$ instead of $H^1$ itself. The
left ($\omega^\pm$) column integrates the \emph{holomorphic} representative
$\Delta$ --- these are Eichler--Shimura's classical periods, your
$\int_0^{i\infty}\Delta\,(X-\tau Y)^{10}d\tau$. The right ($\eta^\pm$)
column integrates the \emph{second-kind} representative $\Dhat$ --- the
quasi-periods. ``Quasi-period'' is the right word for exactly the reason
it is on the curve: it is the period of the second-kind partner of the
period.

\begin{slogan}
$\Delta : \Dhat = \omega : \eta = $ first kind : second kind. The
quasi-periods of $\Dhat$ are the modular analog of the quasi-periods
$\eta_i = \int_{\gamma_i} x\,dx/y$ of an elliptic curve. Same matrix, same
determinant story, one symmetric power up.
\end{slogan}

\section{So where does the modular determinant come from?}

Exactly where Legendre's does, with one distinction worth stating up
front because it is easy to get wrong: the relevant pairing is the
\emph{determinant of the period matrix} $P_f$ --- the comparison
isomorphism between de Rham and Betti realizations --- and this is
\textbf{not} the same as the cup product on cohomology. (The cup product
is a different invariant; see the box below.) The determinant of $P_f$
is the scaling factor by which the comparison isomorphism relates the
rational de Rham structure to the rational Betti structure, and that
scaling is forced to be a rational multiple of the normalizing
$(2\pi i)^{k-1}$. The power is $k-1 = n+1$ because the realizations live
on $\Sym^n$ of $H^1$, not $H^1$.

Numerically (next section), the rational is exactly $n! = 10!$:
\[
  \omega^+\eta^- - \omega^-\eta^+ \;=\; n!\,(2\pi)^{n+1}
  \;=\; 10!\,(2\pi)^{11}\quad(k=12).
\]
In Brown's normalization the $\omega^\pm$ are defined
only up to $K^\times$ (the Hecke field) and the $\eta^\pm$ up to $K^\times$
\emph{and} up to adding $K$-multiples of $\omega^\pm$ --- which is exactly
the statement that the second-kind class is well-defined only modulo the
first-kind one (you can always add a holomorphic form to a second-kind
form). That ``modulo $\omega^\pm$'' freedom is the \emph{same} ambiguity as
your $\CC(X^n - Y^n)$ / $\CC Y^n$ slack: it is the de Rham analog of ``a
second-kind differential is defined up to first-kind differentials and
exact forms.''

\medskip
\noindent\fbox{\parbox{0.95\textwidth}{\small
\textbf{Determinant vs.\ cup product --- do not conflate them (a
correction).} An earlier version of this note claimed the determinant
$\omega^+\eta^- - \omega^-\eta^+$ \emph{is} the cup pairing
$\langle[C^\Delta],[C^{\Dhat}]\rangle$ on $H^1$. That is wrong, and a
direct computation shows why. The invariant symplectic pairing
$\langle\cdot,\cdot\rangle$ on $V_n$ \emph{diagonalizes by parity}: it
couples the even part with the even part and the odd with the odd
(numerically, $\langle B^+, B^-\rangle = 0$ for the Zagier bracket
polynomials). So the cup product of two period polynomials is a
\emph{parity-diagonal} combination,
$\;\omega^+\eta^+\langle B^+,B^+\rangle - \omega^-\eta^-\langle
B^-,B^-\rangle$, which couples $+$ with $+$. The determinant couples $+$
with $-$. They are different bilinear forms in the same four numbers, and
no normalization turns one into the other.
\\[4pt]
What each one computes: the \emph{cup product} (Haberland's formula) is
proportional to the \emph{Petersson norm} --- for $f = g = \Delta$ it
returns $i(\Delta,\Delta)$, matching Zagier's
$\omega^+\omega^-/i(\Delta,\Delta) = 2^{10}$ up to the $V_n$-pairing
normalization. The \emph{determinant} of $P_f$ is the comparison-isomorphism
volume ratio --- the Legendre/quasi-period relation, $=10!\,(2\pi)^{11}$.
Both are real, finite, computable; they are simply different invariants,
one diagonal and one cross in the parity grading.
}}

\medskip
Two cross-checks on the determinant reading, both from your project
knowledge:

\emph{(i) Zagier 1991, $k=12$.} Zagier records
$\omega^+_\Delta \omega^-_\Delta / i(\Delta,\Delta) = 2^{10} \in \QQ$,
where $(\Delta,\Delta)$ is the Petersson norm. This is the
\emph{diagonal} (cup-product / Petersson) statement --- the partner of, but
distinct from, the off-diagonal determinant. It confirms the
rational-times-normalization pattern on the diagonal.

\emph{(ii) Brown 1710.07912 §1.3.} The ``single-valued period'' $s(f)$ is
built as $P^{-1}\overline{P}$ (Fonseca eq.\ 4), and Brown's eq.\ (1.6)
writes $s(f)$ with denominator $\eta^-\omega^+ - \eta^+\omega^- = \det P_f$
(up to sign/$i$), explicitly noting $\det(P_f)$ is proportional to a power
of $2\pi i$ and that $\omega^+\omega^-$ is related to the Petersson norm.
The determinant appears there precisely as the normalizing denominator of
a comparison-isomorphism construction --- i.e.\ as the symplectic scaling
factor.

\section{What this does and doesn't pin down}

If the determinant were a free parameter you could tune, matching it would
mean something. It isn't: $\det P_f = (\text{rational})\cdot(2\pi i)^{k-1}$
is a structural identity (the de Rham/Betti comparison volume ratio), so
\emph{any} four numbers that are genuine periods/quasi-periods satisfy it.
This cuts both ways. If your $\omega^\pm, \eta^\pm$ match Brown's tabulated
values --- which they do, to 20 digits --- then the determinant relation
holds automatically, and ``verifying'' it is checking Brown against Brown
through your arithmetic. It is empty as a test of either Brown or your
construction.

I previously proposed computing the cup pairing
$\langle[C^\Delta],[C^{\Dhat}]\rangle$ from cocycle data as an
\emph{independent} route to the determinant. That proposal was based on a
conflation, now corrected (see the box in \S4): the cup pairing is
parity-diagonal and computes the \emph{Petersson norm}, not the
\emph{cross}-parity determinant. So computing the cup product does not give
an independent handle on $\omega^+\eta^- - \omega^-\eta^+$; it gives
$i(\Delta,\Delta)$ instead. There is no quick cocycle-pairing check of the
determinant, because the determinant simply is not that pairing.

What \emph{would} carry independent information is more pedestrian and you
are already most of the way there: your nine-way agreement across distinct
basepoints $\tau_0$ already stresses the extraction pipeline hard, and the
fact that the \emph{same} finite-contour formula yields $\omega^\pm(\Delta)$
and $\eta^\pm(\Dhat)$ by only changing the integrand is the real structural
content --- that is what places periods and quasi-periods on one footing,
and it is verified by the construction itself, not by any external pairing.

\medskip
\noindent\textbf{The constant, pinned numerically.} From Brown's four
$k=12$ values one finds, to $16$ digits,
\[
  \omega^+\eta^- - \omega^-\eta^+ \;=\; 10!\,(2\pi)^{11}
  \;=\; n!\,(2\pi)^{n+1}
  \qquad (n = 10),
\]
i.e.\ the ``rational $\times$ normalization'' constant is exactly the
factorial $n! = 3628800$ (with $(2\pi)^{n+1}$; equivalently $\pm n!\,(2\pi
i)^{n+1}$, the sign/$i$ depending on the convention carried for the
imaginary $\omega^-,\eta^-$). This used only the four numerical values and
$(2\pi)^{11}$ --- no $V_n$-pairing normalization entered it, so the $n!$ is
solid. It is not a test of Brown (these are his numbers, and the relation
is a theorem they satisfy by construction); it merely records the specific
rational, in case you state the Legendre relation in print.

\bigskip\hrule\bigskip

\noindent\textbf{The note in four sentences.}
An elliptic curve has a rank-$2$ de Rham cohomology with a first-kind
(holomorphic) differential $dx/y$ and a second-kind (residue-free
meromorphic) partner $x\,dx/y$; integrating them over the two torus cycles
gives the period matrix, whose determinant is Legendre's $2\pi i$ because
the matrix intertwines the rational de Rham and Betti intersection
pairings. Brown's theorem says weakly holomorphic modular forms modulo
$D^{n+1}$(lower-weight) \emph{are} the de Rham cohomology of $\Sym^n$ of
that family, with $\Delta$ the first-kind class and $\Dhat$ its second-kind
partner, so $\omega^\pm$ (periods of $\Delta$) and $\eta^\pm$
(quasi-periods of $\Dhat$) are the two columns of the same kind of period
matrix, one symmetric power up. The determinant $\omega^+\eta^- -
\omega^-\eta^+ = n!\,(2\pi)^{n+1}$ (numerically $10!\,(2\pi)^{11}$ at
$k=12$) is the modular Legendre relation: the de Rham/Betti comparison
volume, forced to be rational $\times (2\pi i)^{k-1}$. This is distinct
from the cup product on cohomology, which is parity-diagonal and computes
the Petersson norm instead --- so there is no shortcut ``cup-pairing''
check of the determinant, and matching it against Brown's values is
empty (the relation holds for any genuine periods by construction); the
real structural content of your construction is that one finite-contour
formula yields both $\omega^\pm(\Delta)$ and $\eta^\pm(\Dhat)$.

\end{document}
